This thesis presents the design and implementation of \textit{Favigon}, a full-stack web platform for creating, editing, sharing and exporting web interfaces. The project started from a practical problem that appears frequently in modern product teams: design tools, implementation tools and collaboration tools are often disconnected, which makes the design-to-code process slower and more error-prone than it should be.

The proposed solution combines several capabilities inside a single platform. It includes a visual canvas editor, AI-assisted interface generation, multi-framework code export and a social area where projects can be published, reused and discussed. From a technical perspective, the application uses Angular 20 on the frontend, ASP.NET Core 9 on the backend, PostgreSQL for persistence and Docker Compose for orchestrating the main services. External exposure is handled through Cloudflare Tunnel.

A major part of the thesis focuses on the canvas editor. The editor is treated as a dedicated subsystem with its own state model, history, auto-save behavior, gesture handling, clipboard operations and multi-page support. On the frontend, the logic is separated into focused services responsible for viewport management, persistence, keyboard input, gestures, rendering and code generation. This organization helped keep the implementation manageable while still supporting a large number of interactive behaviors.

The central technical contribution of the work is the AI-assisted design generation and the conversion pipeline from visual design to code. Instead of directly generating final code from a prompt, the platform uses a three-phase pipeline. The first phase extracts the user intent and defines the high-level page blueprint. The second phase creates a structured intermediate representation of the interface. The third phase applies the final visual styling while preserving the validated structure. This staged approach improves controllability and makes error handling more practical.

The same intermediate representation is then reused by the converter engine to generate HTML, React and Angular output. This decision decouples the editor from the target framework and makes the export process easier to extend. The platform also supports responsive export and multi-page generation.

The thesis also covers the surrounding product capabilities required for a usable platform: JWT-based authentication through cookies, OAuth login, password reset, two-factor authentication, project management, public profiles and social actions such as like, star, follow and fork. Security and validation are discussed separately, together with the current automated test suite, which contains 108 passing backend tests.

Overall, the project is not positioned as a replacement for established tools, but as an integrated environment for rapid prototyping, assisted design generation, code export and project sharing. The main personal contributions are the overall system integration, the intermediate UI model, the three-phase AI pipeline and the service-oriented architecture of the canvas editor.
